\documentclass[../../main.tex]{subfiles}
\begin{document}

\begin{center}
  {\fontsize{48pt}{58pt}\selectfont \textbf{1984}\normalsize\textbf{年度}}
\end{center}

\begin{table}[htb]
  \centering
  \caption{出題テーマと難易度}
  \begin{tabular}{|c|p{9cm}|c|}
    \hline
    \textbf{問題番号} & \textbf{テーマ} & \textbf{難易度} \\
    \hline
    第1問 & 円板の影の輪郭と体積 & \\
    \hline
    第2問 & 対数曲線に接する円の中心 & \\
    \hline
    第3問 & 多項式の重根と基底表示 & \\
    \hline
    第4問 & 正三角形の回転体の体積 & \\
    \hline
    第5問 & 細胞の増殖と世代ごとの確率 & \\
    \hline
    第6問 & 放物線領域と点対称な領域の共通部分 & \\
    \hline
  \end{tabular}
\end{table}

\begin{table}[htb]
  \centering
  \caption{解答}
  \renewcommand{\arraystretch}{1.4}
  \begin{tabular}{|c|c|p{9cm}|}
    \hline
    \textbf{問題番号} & \textbf{小問} & \textbf{答え} \\
    \hline
    第1問 & (1) & 影の輪郭は$x^2+(y-\sqrt{2})^2=2$ \\
    \cline{2-3}
    & (2) & $V=\dfrac{2\pi}{3}(2\sqrt{2}-1)$ \\
    \hline
    第2問 & (1) & 極限値は$0$ \\
    \cline{2-3}
    & (2) & 極限値は$1$ \\
    \hline
    第3問 & (1) & $(x-1)^2$で割り切れる条件は所定の表示ができること \\
    \cline{2-3}
    & (2) & $(x-1)^3$で割り切れる条件は係数の和の条件を満たすこと \\
    \hline
    第4問 & -- & $V=\dfrac{\sqrt{3}\pi}{48}$ \\
    \hline
    第5問 & -- & $P_n(1)=p^n$，$P_n(2)=(1-p^n)p^{n-1}$，$P_n(3)=\dfrac{2}{1+p}\left\{p^{3n-2}-(1+p)p^{2n-2}+p^{n-1}\right\}$ \\
    \hline
    第6問 & -- & $0<a<2$，$\max\{a^2,(a-2)^2\}\le b<2(a-1)^2+2$（図示は解答参照） \\
    \hline
  \end{tabular}
\end{table}

\paragraph{試験形態}
6問．試験時間は要確認．

\end{document}
